% STATUS: draft
% PROMOTE: exemplars
% TAGS: arrays, math
% DIFFICULTY: easy
% SOURCE: LeetCode 66

\ProblemTitle{Plus One}
\ProblemMeta{Easy}{Arrays, Simulation}{LC-066}
\label{prob:lc-066-plus-one}

\ProblemSummary
Given a non-empty array of digits representing a non-negative integer, increment
the integer by one and return the resulting array of digits.

The most significant digit is stored at the beginning of the array, and each
digit is in the range \([0,9]\).

\KeyIdea
Incrementing the number requires simulating elementary addition from right to
left. Carry propagation only affects trailing digits equal to nine; once a
digit strictly less than nine is incremented, the process terminates.

\begin{Thinking}
My initial instinct was to convert the digit array into an integer, add one, and
convert it back. This approach fails due to potential overflow and violates the
problem constraints.

What simplified the reasoning was treating the operation as grade-school
addition and maintaining the invariant that all digits to the right of the
current index are already correct after carry handling.
\end{Thinking}

\Algorithm
\begin{CodeBlock}[style=pseudo,caption={Pseudocode}]{12}
for i from n - 1 downto 0:
    if digits[i] < 9:
        digits[i] = digits[i] + 1
        return digits
    digits[i] = 0

prepend 1 to digits
return digits
\end{CodeBlock}

\Correctness
We iterate from the least significant digit toward the most significant digit.
At each iteration, the invariant maintained is that all digits to the right of
the current index represent the correct result after increment and carry
propagation.

If a digit less than nine is encountered, incrementing it produces the correct
result and no further carry is required. If all digits are nine, each is set to
zero and a leading one is prepended, yielding the correct incremented number.
Thus, the algorithm always returns the correct result.

\Complexity
The algorithm runs in \(O(n)\) time, where \(n\) is the number of digits. It uses
\(O(1)\) additional space beyond the output array.

\bigskip
\noindent\textbf{C++ Implementation}

\lstinputlisting[
  style=cpp,
  caption={C++ Solution — Plus One}
]{../../../../problems/01-arrays-sequences/easy/lc-066-plus-one/solution.cpp}

\bigskip
\noindent\textbf{Python Implementation}

\lstinputlisting[
  language=Python,
  backgroundcolor=\color{codebg},
  basicstyle=\ttfamily\small,
  frame=single,
  breaklines=true,
  caption={Python Solution — Plus One}
]{../../../../problems/01-arrays-sequences/easy/lc-066-plus-one/solution.py}

\ProblemDivider
